\documentclass[letterpaper,10pt]{article}
\usepackage[margin=0.5in]{geometry}
\usepackage[T1]{fontenc}
\usepackage[utf8]{inputenc}
\usepackage{helvet}
\renewcommand{\familydefault}{\sfdefault}
\usepackage{enumitem}
\usepackage[hidelinks]{hyperref}
\setlength{\parindent}{0pt}
\setlength{\parskip}{2pt}
\pagestyle{plain}
\begin{document}
{\LARGE\textbf{Alexander Tong}}\\
860-706-3624 | alto335@bu.edu | alexandertong.space | linkedin.com/in/alto335\\
\section*{Education}
Boston University, Boston, MA | Expected May 2028\\
B.S. Mechanical Engineering, Aerospace Engineering Concentration | GPA: 3.71/4.00\\
73 earned semester hours of BU coursework; 24 AP test credits listed separately on transcript\\
Completed: Energy and Thermodynamics; Mechanics of Materials; CAD and Machine Components\\
In progress: Fluid Mechanics; Manufacturing Processes; Electromechanical Design\\
\section*{Engineering Experience and Projects}
\textbf{Argo Separation and Recovery Subsystem}\hfill 06/2025 - Present\\
Co-Lead Engineer | Boston University Rocket Propulsion Group | Boston, MA\\
Student club project | Total BURPG commitment: 10-30 hours/week, variable
\begin{itemize}[leftmargin=12pt,itemsep=2pt,topsep=2pt,parsep=0pt]
\item Designed redundant recovery for a 135 lb liquid rocket targeting 100,000 ft; achieved 18.7 lbm against a 20 lbm budget and presented four iterative design reviews
\item Compared system architectures, commercial flight computers, and separation mechanisms against mass, cost, and development constraints; recovery system cost was \$1,076
\item Calculated 823 lbf parachute shock and 1,230 lbf snatch loads by synthesizing five research-paper methods; introduced and dynamically tested Brummel eye splices to replace knots and reduce quick-links, confirming minimum 1.5 FOS
\item Developed CO2 deployment hardware after recurring integrated-test failures and parachute damage; contributed pressure and Lamé thick-wall stress calculations, Fusion 360 CAD, drawings, machining, and hardware tests in a three-person team; integrated CO2 validation remains planned
\item Led recovery test operations: test plans, safety protocols, crew assignments, briefings, firing authorization, abort decisions, and post-test analysis
\item Built and operated a 352 L vacuum chamber at 0.011 atm using salvaged pumps and available stock; integrated LabJack data acquisition, a pressure transducer, and two-amplifier signal conditioning for deployment-hardware and avionics tests
\item Redesigned and machined cable cutters on manual mill and lathe, reducing mass 49.5\% with a calculated 2.01 FOS; documented assembly and recovery-system packing procedures
\end{itemize}
\textbf{Boston University Plummer Laboratory}\hfill 06/2026 - Present\\
Mechanical Engineering Intern | Boston, MA\\
06/2026 - 08/2026: full-time, 40 hours/week | 09/2026 - Present: part-time, 10 hours/week
\begin{itemize}[leftmargin=12pt,itemsep=2pt,topsep=2pt,parsep=0pt]
\item CNC-machined a 6061-T6 shaker adapter on a Haas VF-2; reduced stock mass 40.4\%, from 3.19 to 1.90 kg, with a conical underside chamfer, corner fillets, and 12 drilled holes
\item Generated G-code, engineering drawings, and tolerance stack-ups for +/-0.005 in holes; reduced estimated plate-and-fastener cost from \$1,170 outsourced to \$71.11 through in-house fabrication and fastener selection
\item Predicted 0.44 mm fixture deflection under a 3 kg load at 20.1 Hz using SolidWorks dynamic finite element analysis (FEA); analyzed fastener preload and joint failure modes
\item Independently measured 50-120 mN forces with an Instron 68SC-5 to investigate physical Ising models and support ongoing comparison with mathematical predictions
\item Presented shaker-adapter design iterations, finished hardware, and results in a flash-talk update at a meeting involving multiple laboratories
\end{itemize}
\newpage
\section*{Engineering Experience and Projects Continued}
\textbf{Culturon}\hfill 02/2026 - 06/2026\\
Mechanical Engineering Intern | Sydney, Australia\\
Part-time | 10 hours/week
\begin{itemize}[leftmargin=12pt,itemsep=2pt,topsep=2pt,parsep=0pt]
\item Planned and executed the company's first full-scale nanoparticle deposition mapping across chamber locations and magnetic configurations; used scanning electron microscopy (SEM), optical microscopy, and ImageJ to characterize polarity-dependent deposition and asymmetric flow behavior
\item Formed a working hypothesis for contamination cause and pattern from sample-plate SEM/ImageJ results, optical microscopy, and direct visual evidence
\item Presented turbomolecular- and roughing-pump port, chamber-bottom, and electron-curtain prevention concepts based on the hypothesis and fluid/electromagnetic first principles in a technical design review; full-system implementation remained conceptual
\item Designed Onshape sheet-metal sample racks and coupons for reactor-port insertion and adhesive-free mounting; incorporated bend reliefs, fabrication tolerances, assembly clearances, and flow-disturbance considerations
\item Compared observations with published research and presented findings and design recommendations to technical teams; proprietary results and parameters are protected by NDA
\end{itemize}
\textbf{AIAA Design-Build-Fly Aircraft}\hfill 06/2025 - 12/2025\\
Wing Structural Lead | Boston University | Boston, MA\\
Student club project | 10-20 hours/week, variable
\begin{itemize}[leftmargin=12pt,itemsep=2pt,topsep=2pt,parsep=0pt]
\item Led two wing designs in Fusion 360; compared structure, material, and manufacturing options to produce a 50\% lighter XPS prototype and reduce build time from three months to one
\item Designed a split NACA 2412 wing with 0.52 taper ratio and 5-degree dihedral using aerodynamic-team specifications and manufacturing lessons from the first prototype; developed the servo flap system
\item Led structural and manufacturing reviews, mentored nine members in laser cutting, waterjet operation, and Fusion 360, and coached the succeeding wing lead in manufacturing logistics
\end{itemize}
\textbf{Hybrid Rocket Nozzle}\hfill 09/2024 - 11/2024\\
Lead Engineer | Boston University Rocket Propulsion Group | Boston, MA\\
Part-time student club project | Approximately 5 hours/week
\begin{itemize}[leftmargin=12pt,itemsep=2pt,topsep=2pt,parsep=0pt]
\item Led five members through three design reviews to develop a parabolic nozzle using isentropic flow relations, Rao contour methods, and SolidWorks CAD/FEA
\item Reduced nozzle mass 73\% from the initial design and calculated a 7.8 FOS under mechanical loads alone; incorporated review feedback on manufacturing and assembly fragility
\item Led the team developing the launch computer and completed an approximately two-second hotfire; a load-cell malfunction prevented thrust-data collection
\end{itemize}
\textbf{Ablative Anode Pulsed Plasma Thruster Laboratory}\hfill 05/2026 - 06/2026\\
Volunteer | University of Sydney | Sydney, Australia\\
Part-time volunteer work | Approximately 5 hours/week
\begin{itemize}[leftmargin=12pt,itemsep=2pt,topsep=2pt,parsep=0pt]
\item Designed an ABS inductor spool, performed coil-winding calculations, and checked vacuum-outgassing suitability; provided limited assistance with the established thruster operating procedure
\end{itemize}
\textbf{Space Physics and Technology Laboratory}\hfill 10/2025 - 12/2025\\
COSSMo Research Experience | Boston University | Boston, MA\\
Part-time student research | Approximately 5 hours/week
\begin{itemize}[leftmargin=12pt,itemsep=2pt,topsep=2pt,parsep=0pt]
\item Assisted with and observed flight-spare laser calibration; combined existing research tools, pySPEDAS, and self-written Python functions for limited analysis of live and historical spacecraft data
\end{itemize}
\section*{Technical Skills}
Design and analysis: SolidWorks, Fusion 360, Onshape, FEA, GD\&T, engineering drawings, MATLAB, Python\\
Fabrication and testing: CNC and manual machining, composite layup, 3D printing, vacuum systems, LabJack data acquisition, pressure instrumentation, test planning and operations\\
\end{document}
